% ============================================================================
%  SECCIÓN 1.1: ANTECEDENTES
% ============================================================================

\section{Antecedentes}

Los antecedentes consisten en la presentaci\'on de la informaci\'on m\'as
relevante y directamente relacionada con el tema del proyecto. Su estudio
garantiza que el trabajo parta del punto m\'as actual de la tem\'atica en
cuesti\'on. Para su elaboraci\'on se emplean los m\'etodos hist\'orico l\'ogico y
deductivo, partiendo de elementos generales, como el desarrollo de la
infraestructura deportiva del pa\'is, hasta llegar a los aspectos espec\'ificos
relacionados con la difusi\'on de la oferta deportiva y el uso de tecnolog\'ias de
informaci\'on en este \'ambito.

\subsection{Antecedentes hist\'oricos}

El deporte y su infraestructura han sido parte de la vida social del pa\'is
desde comienzos del siglo XX. Escenarios emblem\'aticos como el Estadio Hernando
Siles de La Paz, inaugurado en 1930, el Estadio F\'elix Capriles de
Cochabamba y el Estadio Ram\'on Tahuichi Aguilera de Santa Cruz de la Sierra
han sido, durante d\'ecadas, los centros de la pr\'actica deportiva y del
encuentro social de las comunidades. Con el paso del tiempo, la oferta de
recintos se ha diversificado: estadios, coliseos cerrados, canchas de
baloncesto y voleibol, frontones, piscinas, pistas de atletismo y complejos
polideportivos distribuidos en los nueve departamentos del pa\'is.

Sin embargo, la informaci\'on sobre estos espacios siempre ha estado
fragmentada y ha dependido de la difusi\'on que cada recinto realiza de manera
individual mediante medios impresos, radiales, televisivos o redes sociales.
Esta situaci\'on hist\'orica explica que, a pesar de contar con una
infraestructura deportiva considerable, la poblaci\'on no disponga de un medio
unificado para conocerla.

\subsection{Antecedentes tem\'aticos}

A nivel internacional existen plataformas consolidadas que integran
informaci\'on georreferenciada sobre espacios deportivos, de recreaci\'on y de
esparcimiento, tales como directorios de instalaciones deportivas, sistemas
de reserva de canchas y aplicaciones de mapas que combinan cat\'alogo,
fotograf\'ias, horarios y eventos. Estas soluciones han demostrado que la
centralizaci\'on de la informaci\'on incrementa la visibilidad de los recintos,
favorece el uso de la infraestructura existente y facilita la toma de
decisiones de los usuarios \citep{longley2015, pressman2010}.

No obstante, la mayor\'ia de estas plataformas responden a contextos de
pa\'ises desarrollados y no se adaptan a las particularidades del medio
nacional, donde la informaci\'on deportiva es escasa y poco estructurada. El
tema ha sido abordado de forma limitada en la literatura local, y los pocos
intentos de digitalizaci\'on se concentran en recintos individuales o en
instituciones espec\'ificas, sin una visi\'on integradora a nivel pa\'is.

\subsection{Antecedentes institucionales}

El proyecto se enmarca en el \'ambito de las instituciones deportivas y
educativas del pa\'is. Los escenarios deportivos, tanto p\'ublicos como
privados, as\'i como las federaciones deportivas, los clubes y las
direcciones municipales de deportes, constituyen la poblaci\'on institucional
objetivo del sistema. Estas instituciones comparten la necesidad com\'un de
visibilizar su oferta deportiva, pero carecen de herramientas tecnol\'ogicas
adecuadas para hacerlo, lo que justifica la propuesta de una plataforma
centralizada.

\subsection{Investigaciones previas}

Diversos estudios en el \'ambito de los sistemas de informaci\'on geogr\'afica
(SIG) aplicados a servicios urbanos han demostrado que la visualizaci\'on
cartogr\'afica de la infraestructura mejora la accesibilidad a la informaci\'on
y apoya la planificaci\'on de pol\'iticas p\'ublicas \citep{longley2015}.
Trabajos similares, orientados al turismo y a la recreaci\'on, han utilizado
mapas interactivos y bases de datos centralizadas con resultados favorables
en t\'erminos de difusi\'on y participaci\'on de los usuarios. Estos hallazgos
orientan la metodolog\'ia y el dise\~no de la presente propuesta.

\subsection{Estad\'isticas consultadas}

Las fuentes estad\'isticas consultadas reflejan una baja participaci\'on de la
poblaci\'on en la actividad f\'isica y deportiva, atribuida en gran medida a la
falta de informaci\'on sobre la oferta disponible. Encuestas de h\'abitos
deportivos se\~nalan que el desconocimiento de la ubicaci\'on de los escenarios,
de las disciplinas que albergan y de los horarios de uso es una de las
principales barreras para la pr\'actica deportiva. Estas cifras respaldan la
pertinencia de la propuesta y sirven de base para la definici\'on de los
indicadores del proyecto \citep{hernandez2014}.

\begin{figure}[H]
  \centering
  \includegraphics[width=0.62\linewidth]{escenarios-deportivos.jpeg}
  \caption{Imagen ilustrativa de los escenarios deportivos del pa\'is}
  \label{fig:infraestructura-deportiva}
\end{figure}

\textbf{FUENTE:} Elaboraci\'on propia en base a la revisi\'on documental del
sector deportivo.
